\documentclass[11pt]{article}
\usepackage[T1]{fontenc}
\usepackage{lmodern,amsmath,amssymb,amsthm,booktabs,enumitem}
\usepackage[margin=1in]{geometry}
\usepackage[hidelinks]{hyperref}
\newtheorem{theorem}{Theorem}[section]
\newtheorem{proposition}[theorem]{Proposition}
\newtheorem{lemma}[theorem]{Lemma}
\newtheorem{corollary}[theorem]{Corollary}
\theoremstyle{remark}\newtheorem{remark}[theorem]{Remark}
\newcommand{\gen}[1]{\langle #1\rangle}
\newcommand{\ncl}[1]{\langle\!\langle #1\rangle\!\rangle}
\newcommand{\Ctwo}{C_2}
\newcommand{\IP}{\mathrm{IP}}
\title{Finite quotients with prescribed parabolic intersections}
\author{human\_compression research project}
\date{11 October 2026\quad---\quad Phase 281, local mathematical note}
\begin{document}
\maketitle
\begin{abstract}
We prove two sufficient conditions for finite quotients to preserve prescribed
subgroup intersections. The first uses endomorphisms of a normal subgroup of
finite index. The second applies to a marked string group with one finite
maximal parabolic, and replaces simultaneous control of all parabolics by a
subgroup-separability hypothesis on the opposite parabolic. We also give the
precise consequence of uniqueness of a finite minimal regular cover for its
two end parabolics, and a finite-index obstruction when a proposed fixed facet
is flat. These results do not settle the remaining all-rank question about
finite minimal regular covers. The unproved filling hypothesis is stated
separately from the propositions proved here.
\end{abstract}

\section{A normal finite quotient from endomorphisms}

All subgroups in this section are ordinary abstract subgroups. An index set
\(\mathcal I\) is finite and closed under an operation denoted by intersection;
in the intended application its members are subsets of the set of marks.
Write \(P_{I\cap J}\) for the specified common subgroup.

\begin{theorem}\label{thm:normal}
Let \(T\) be finitely generated, and let \((P_I)_{I\in\mathcal I}\) be a finite
family of subgroups such that
\[
 P_I\cap P_J=P_{I\cap J}.
\]
Suppose \(L\lhd T\) has finite index and there are endomorphisms
\(\rho_I:L\to L\) satisfying
\begin{align}
 \rho_I(a)&=a &&(a\in P_I\cap L),\label{eq:fix}\\
 \rho_I(P_J\cap L)&\subseteq P_{I\cap J}\cap L
      &&(I,J\in\mathcal I).\label{eq:cross}
\end{align}
Suppose also that some normal subgroup \(D\lhd T\) of finite index satisfies
\begin{equation}\label{eq:approx}
 D\subseteq L,\qquad
 P_ID\cap P_JD\subseteq P_{I\cap J}L
       \quad(I,J\in\mathcal I).
\end{equation}
Then for every normal subgroup \(M\lhd T\) of finite index there is a normal
subgroup \(N\lhd T\) of finite index, with \(N\subseteq M\), such that
\begin{equation}\label{eq:exact}
 P_IN\cap P_JN=P_{I\cap J}N\quad(I,J\in\mathcal I).
\end{equation}
The endomorphisms \(\rho_I\) need not commute.
\end{theorem}

\begin{proof}
Put \(D_0=D\cap M\) and \(F=L/D_0\). The group \(L\) is finitely generated
because it has finite index in \(T\). Thus there are only finitely many
homomorphisms \(L\to F\): each is determined by its values on a fixed finite
generating set, and \(F\) is finite. Define
\[
 N=\bigcap_{f\in\operatorname{Hom}(L,F)}\ker f.
\]
Every kernel in this intersection has finite index, so \(N\) has finite
index in \(L\), and hence in \(T\). The natural quotient \(L\to F\) occurs
among these homomorphisms. Consequently \(N\subseteq D_0\subseteq M\).

If \(\alpha:L\to L\) is any endomorphism and \(x\in N\), then
\(f(\alpha(x))=(f\circ\alpha)(x)=1\) for every \(f:L\to F\). Hence
\(\alpha(N)\subseteq N\). In particular every \(\rho_I\) preserves \(N\).
Conjugation by an element of \(T\) induces an automorphism of \(L\), since
\(L\lhd T\). Applying the same observation to this automorphism and its
inverse proves \(N\lhd T\).

It remains to prove the nontrivial inclusion in \eqref{eq:exact}. Choose
\(a\in P_I\) and \(b\in P_J\) with \(aN=bN\). Since \(N\subseteq D\),
condition \eqref{eq:approx} gives \(a\in P_{I\cap J}L\). Write \(a=ca_0\),
where \(c\in P_{I\cap J}\) and \(a_0\in L\). Then \(a_0\in P_I\cap L\).
Put \(b_0=c^{-1}b\). From \(aN=bN\) and normality we obtain
\(a_0b_0^{-1}\in N\subseteq L\), so \(b_0\in P_J\cap L\). Applying
\(\rho_I\), and using \eqref{eq:fix}--\eqref{eq:cross}, yields
\[
 a_0\rho_I(b_0)^{-1}\in N,\qquad
 \rho_I(b_0)\in P_{I\cap J}.
\]
It follows that \(aN=c\rho_I(b_0)N\) belongs to the image of
\(P_{I\cap J}\). The reverse inclusion is immediate.
\end{proof}

The right-hand side of \eqref{eq:approx} contains \(L\), not \(D\).
Thus the hypothesis is weaker than the desired intersection equality in a
finite quotient. The additional condition is a compatible family of
endomorphisms on one \emph{normal} subgroup of finite index. This normality
is part of the theorem, not something obtained merely by replacing an
arbitrary finite-index subgroup by its core.

\section{A finite facet and a separable opposite parabolic}

A subgroup \(H\leq T\) is \emph{separable} if, for every \(x\notin H\),
there is a homomorphism \(q\) from \(T\) to a finite group such that
\(q(x)\notin q(H)\). No normality of \(H\) is assumed.

We use the standard intersection and quotient criteria for string C-groups.
For an ordered string of involutions, the intersection criterion says that
the two end parabolics being C-groups, with intersection equal to the common
middle parabolic, implies the full intersection property. The quotient
criterion says that a marked epimorphism onto a string C-group, injective
on one end parabolic, makes the source a string C-group. These are classical
inputs, also used in the proof of Monson--Schulte's Lemma 3.3
\cite{MS}. They are not consequences of separability. To recall their
relationship, one proves the quotient criterion by induction on rank:
the opposite end satisfies the same quotient criterion one rank lower,
while injectivity on the first end converts the intersection of their
images into equality with the middle parabolic. The intersection criterion
then completes the induction.

\begin{theorem}\label{thm:facet}
Let \(T=\gen{s_0,\ldots,s_{n-1}}\), with \(n\geq3\), be a group generated
by distinct nonidentity involutions satisfying the distant commutations.
Put
\[
 A=\gen{s_0,\ldots,s_{n-2}},\quad
 B=\gen{s_1,\ldots,s_{n-1}},\quad
 C=\gen{s_1,\ldots,s_{n-2}}.
\]
Assume:
\begin{enumerate}[label=(\roman*),leftmargin=2em]
\item \(A\) is a finite string C-group, and \(A\cap B=C\);
\item there is a marked epimorphism \(\beta:B\to B_0\), where \(B_0\)
is a finite string C-group and \(\beta\) is injective on \(C\);
\item \(T\) is residually finite and \(\ker\beta\) is separable in \(T\).
\end{enumerate}
For every homomorphism \(\theta:T\to G\) to a finite group, there is a
finite marked quotient \(q:T\to H\) such that \(H\) is a string C-group,
\(q\) is injective on \(A\), and \(\theta\) factors through \(q\).
In particular, all relations already holding in \(T\) still hold in \(H\).
\end{theorem}

\begin{proof}
Write \(L=\ker\beta\). Since \([B:L]\) is finite, \(B\) is a finite union
of cosets of the separable subgroup \(L\), and is therefore separable in
\(T\). Here is the finite-quotient argument explicitly. Choose representatives
\(b_1,\ldots,b_k\) with \(B=\bigcup_i Lb_i\). If \(x\notin B\), then
\(xb_i^{-1}\notin L\) for every \(i\). Separability gives normal subgroups
\(N_i\lhd T\) of finite index such that
\(xb_i^{-1}\notin LN_i\). For \(N_x=\bigcap_iN_i\) this implies
\(x\notin BN_x\). This is the asserted separability of \(B\).

We now choose three finite collections of normal subgroups of finite index.
First, residual finiteness excludes every member of \(A\setminus\{1\}\)
from a normal subgroup \(N_A\) of finite index. Second, for every
\(a\in A\setminus C\), the equality \(A\cap B=C\) gives \(a\notin B\).
Separability of \(B\) supplies a normal subgroup \(N_a\) of finite index
such that \(a\notin BN_a\). There are only finitely many such \(a\).
Third, for each representative \(b_i\notin L\), separability of \(L\)
supplies a normal subgroup \(N'_i\) of finite index with
\(b_i\notin LN'_i\). Intersect all these subgroups, and also \(\ker\theta\),
to obtain a normal subgroup \(N\lhd T\) of finite index.

By construction,
\begin{equation}\label{eq:three}
 N\cap A=1,\qquad A\cap BN=C,\qquad N\cap B\subseteq L.
\end{equation}
For the last inclusion, write \(b\in N\cap B\) as \(lb_i\), with \(l\in L\).
If \(b_i\notin L\), then \(b_i=l^{-1}b\in LN\subseteq LN'_i\), a
contradiction. The other two assertions follow directly from the first and
second collections.

Set \(H=T/N\). Its first end is isomorphic to \(A\). Its opposite end is
\(\overline B=B/(B\cap N)\), which maps onto \(B_0\) because
\(B\cap N\subseteq L\). The common middle subgroup \(\overline C\) is
isomorphic to \(C\), and its map to \(B_0\) is injective. The quotient
criterion applied to \(\overline B\to B_0\) therefore proves that
\(\overline B\) is a string C-group. The second equality in
\eqref{eq:three} gives
\(\overline A\cap\overline B=\overline C\). The intersection criterion
now proves that \(H\) is a string C-group. Finally \(N\subseteq\ker\theta\)
gives the required factorization.
\end{proof}

\begin{remark}
The theorem can preserve a relation such as \(w^m=1\), provided it is
already imposed in \(T\) and all the stated hypotheses hold for that
\(T\). It does not claim that residual finiteness alone preserves
parabolic intersections. The separability of \(\ker\beta\) is essential
to the argument just given.
\end{remark}

\section{What uniqueness would force, and the flat case}

A finite minimal regular cover means a finite regular cover which has no
proper regular quotient still covering the given marked group or polytope.
Different covering maps to the same object are not counted as different
minimal covers. All maps below preserve the distinguished generators.

\begin{lemma}\label{lem:faithful}
Let \(G\) be a finite marked string group whose two end parabolics \(G_A\)
and \(G_B\) are string C-groups. If \(G\) has a unique finite minimal
C-group cover \(M\), then the maps \(M_A\to G_A\) and \(M_B\to G_B\)
are isomorphisms.
\end{lemma}

\begin{proof}
The exact interval-extension argument in phase 277, obtained from the
construction in Monson--Schulte's Lemmas 3.1 and 3.3 \cite{MS}, gives a
finite C-group cover \(R\to G\) whose first end is exactly \(G_A\).
Choose a minimal C-group quotient \(M'\) of \(R\) still covering \(G\).
Such a quotient exists because \(R\) is finite. On first ends the composite
\[
 G_A=R_A\longrightarrow M'_A\longrightarrow G_A
\]
is the identity on every distinguished generator, and hence on the group.
Both arrows are surjective, so both are isomorphisms. Uniqueness identifies
\(M'\) with \(M\). Repeat the argument starting with the other end.
This uses a separately chosen exact extension for each end; it does not
assert that the restriction of an arbitrary minimal cover is minimal.
\end{proof}

\begin{proposition}\label{prop:flat}
Let \(A=\gen{s_0,\ldots,s_{n-2}}\) be a finite string C-group, with
\(n\geq3\), and put
\[
 D=\gen{s_0,\ldots,s_{n-3}},\quad
 C=\gen{s_1,\ldots,s_{n-2}},\quad E=C\cap D.
\]
Let
\[
 U=A*_{D}(D\times\gen{t}),\qquad |t|=2,
 \qquad \widetilde B=\gen{C,t}\leq U.
\]
Then \(\widetilde B\cong C*_{E}(E\times\gen t)\) and
\(A\cap\widetilde B=C\). If \(A=CD\), then
\[
 U=A\widetilde B,\qquad [U:\widetilde B]=[A:C]<\infty.
\]
Consequently every element of \(U\) has a positive power in
\(\widetilde B\).
\end{proposition}

\begin{proof}
The homomorphism from the displayed amalgam for \(\widetilde B\) to \(U\)
is injective by the reduced-word theorem for amalgamated products. Indeed,
elements of \(C\setminus E\) lie in \(A\setminus D\), and elements of
\((E\times\gen t)\setminus E\) lie in
\((D\times\gen t)\setminus D\). Thus a reduced alternating word for the
smaller amalgam stays reduced in the larger one. The same normal form
shows that its intersection with \(A\) is precisely \(C\).

Suppose \(A=CD\). Taking inverses also gives \(A=DC\). The set
\(A\widetilde B\) contains the identity and is preserved by left
multiplication by \(A\). It is also preserved by left multiplication by
\(t\): if \(a=dc\), with \(d\in D\) and \(c\in C\), then
\(tab=dtcb\in A\widetilde B\) for \(b\in\widetilde B\), since \(t\)
commutes with \(D\). As \(A\) and \(t\) generate \(U\), we obtain
\(U=A\widetilde B\). The cosets of \(\widetilde B\) are therefore indexed
by the cosets of \(C=A\cap\widetilde B\) in \(A\). The final assertion
follows from the action of a cyclic subgroup on this finite coset set.
\end{proof}

For the regular polytope with group \(A\), the equality \(A=CD\) says
that every vertex is incident with every facet: the polytope is flat.
Proposition~\ref{prop:flat} therefore prevents an argument which requires
an element having no positive power in the opposite parabolic, while
keeping this particular flat facet unchanged.

\section{The remaining filling hypothesis}

Here is the proposed use of Theorem~\ref{thm:facet}, with its unproved
part made explicit. Suppose the two end parabolics of a finite marked
group \(G\) are C-groups and a unique finite minimal cover \(M\) exists.
Lemma~\ref{lem:faithful} supplies the map \(U\to M\to G\), where \(U\)
is the amalgam in Proposition~\ref{prop:flat}, with \(A=G_A\).
To obtain a contradiction by a power relation, one would need an element
\(w\in\ker(U\to G)\setminus\ker(U\to M)\) and an integer \(m\) coprime
to the order of its image in \(M\), such that
\[
 T=U/\ncl{w^m}
\]
satisfies Theorem~\ref{thm:facet}. In particular the images of \(A\) and
\(\widetilde B\) must have the specified intersection, and the kernel of
a suitable finite C-group quotient of the latter must be separable in
\(T\). The resulting finite C-group cover \(H\to G\) would kill
\(w^m\), whereas its unique minimal descendant \(M\) would not. This
last contradiction is valid only after the stated properties of \(T\)
have been proved.

The present note does not prove those filling properties. In the
nonflat case a possible route passes through a free normal subgroup of
finite index and small cancellation; preservation of the opposite
parabolic and the required primary filling statements have not been
completed here. In the flat case Proposition~\ref{prop:flat} rules out
the proposed choice avoiding powers in \(\widetilde B\). Replacing
\(A\) by a larger finite C-group cover may remove flatness, but the
kernel of the resulting map to \(G\) can then contain torsion. An
argument for free kernels cannot be reused without proving a new
filling statement for that situation. A theorem allowing only
exponents divisible by a fixed integer is also insufficient when that
integer has a common factor with the order of the image of \(w\) in \(M\).

The unresolved question here is the finite-cover strengthening, in
arbitrary rank. The earlier result allowing infinite regular covers,
the finite rank-four theorem, and the interval propagation theorem
remain separate results. No new complete historical manuscript is
claimed to follow from this note, and no new compression factor or
solution of the all-rank finite question is assigned to it.

\section{Attribution and formal verification}

The subgroup-separation and finite-index arguments in
Theorems~\ref{thm:normal} and~\ref{thm:facet} are proved above. They use
standard group-theoretic constructions; this note does not assert
historical priority for those constructions. The string C-group
intersection and quotient criteria are retained classical inputs.
Lemma~\ref{lem:faithful} uses the exact interval-extension proof already
recorded in phase 277, with the Monson--Schulte construction explicitly
retained. Proposition~\ref{prop:flat} uses the classical normal form for
amalgamated free products. The proposed small-cancellation application
in the preceding section is not one of the proved theorems.

The formal file, \texttt{FiniteQuotientIntersections.lean}, verifies six
declarations using Lean 4.32.2 and the pinned mathlib revision recorded
in the replay script. They establish membership and inclusion for the
intersection of all homomorphism kernels, its invariance under every
endomorphism, equality of two subgroup images under an explicit
membership-detection hypothesis, the corresponding quotient statement,
and injectivity on a subgroup under kernel avoidance. These are actual
group, homomorphism, and subgroup statements. The existence of the
finite-index subgroups, the separability assumptions, the classical
C-group criteria, and the amalgam argument remain written proofs or
explicit retained inputs. In particular the Lean declarations do not
formalize a filling theorem or a solution of the remaining question.

\begin{thebibliography}{9}
\bibitem{MS}
B. Monson and E. Schulte,
\emph{Finite Polytopes have Finite Regular Covers},
arXiv:1209.1167, consulted seven-page author PDF with v1 header and
internal title-page date 21 November 2021, Lemmas 3.1 and 3.3,
pp.~4--5. The exact source bytes and earlier phase 277 proof are bound
in the accompanying source record.
\end{thebibliography}
\end{document}
